\section{\zhen{全部结果速查}{All results at a glance}}\label{sec:overview}

\zhen{这一节只供查找，不含新内容。}{This section is for looking things up; it contains nothing new.}

% Experiment-driven papers: replace "All main results" with a method-components table
% (component | role | formula or algorithm | where in the explainer) and delete "Relations between results".
\subsection{\zhen{所有主要结果}{All main results}}

% The statement column uses the corrected form (as in E2 / 附 2). For a conclusion that is the explainer's own,
% write "(explainer's note)" / "（讲解注）" in the first column.
\begingroup
\small
\renewcommand{\arraystretch}{1.35}
\begin{longtable}{@{}>{\raggedright\arraybackslash}p{0.13\linewidth}>{\raggedright\arraybackslash}p{0.5\linewidth}>{\raggedright\arraybackslash}p{\dimexpr0.37\linewidth-4\tabcolsep\relax}@{}}
\toprule
\textbf{\zhen{结果}{Result}} & \textbf{\zhen{内容}{Statement}} & \textbf{\zhen{含义}{Meaning}} \\
\midrule
\endfirsthead
\toprule
\textbf{\zhen{结果}{Result}} & \textbf{\zhen{内容}{Statement}} & \textbf{\zhen{含义（续）}{Meaning (cont.)}} \\
\midrule
\endhead
\bottomrule
\endlastfoot
\FILL{one row per result} & & \\
\end{longtable}
\endgroup

\subsection{\zhen{结果之间的关系}{Relations between results}}

% Keep only when there are many results (6 or more) with non-trivial dependencies; otherwise delete this subsection.
% TikZ dependency graph: arrows from the lemma used to the theorem that uses it; Part I and Part II results separated by a dashed line.
\FILL{dependency graph}

\subsection{\zhen{符号速查}{Notation}}

\begingroup
\small
\renewcommand{\arraystretch}{1.3}
\begin{longtable}{@{}>{\raggedright\arraybackslash}p{0.3\linewidth}>{\raggedright\arraybackslash}p{\dimexpr0.7\linewidth-2\tabcolsep\relax}@{}}
\toprule
\textbf{\zhen{符号}{Symbol}} & \textbf{\zhen{含义（首次出现）}{Meaning (first appearance)}} \\
\midrule
\endfirsthead
\toprule
\textbf{\zhen{符号}{Symbol}} & \textbf{\zhen{含义（首次出现）（续）}{Meaning (first appearance, cont.)}} \\
\midrule
\endhead
\bottomrule
\endlastfoot
\FILL{one row per symbol} & \\
\end{longtable}
\endgroup
